% Part: second-order-logic
% Chapter: syntax-and-semantics
% Section: language-of-sol

\documentclass[../../../include/open-logic-section]{subfiles}

\begin{document}

\olfileid{sol}{syn}{lan}

\olsection{د دويمې درجې منطق ژبه}

لکه په لومړۍ درجه منطق کښې، د دويمې درجې منطق عبارتونه له داسې بنسټيزې زېرمې جوړېږي چې \emph{!!{variable}s}، \emph{!!{constant}s}، \emph{!!{predicate}s} او کله کله \emph{!!{function}s} لري۔ له دې سمبولونو، منطقي تړنونو، کميت‌ښودونکو او د قوسونو او کامو په څېر د ليک نښو څخه \emph{ترمونه} او \emph{!!{formula}s} جوړېږي۔ فرق دا دے چې د څيزونو پر متغيرونو سربېره په دويمه درجه منطق کښې د اړيکو او تابعو متغيرونه هم شته او کميت‌ښودنه پرې کېدے شي۔ نو د لومړۍ درجې منطقي سمبولونه شته، او ورسره دا هم شته:

\begin{enumerate}
\item د اړيکې د ځايونو د هر شمېر~$n$ دپاره د اړيکو د دويمې درجې !!{variable}s يو !!{denumerable}s سټ: $\Obj V_0^n$، $\Obj V_1^n$، $\Obj V_2^n$، \dots
\item د تابعو د دويمې درجې !!{variable}s يو !!{denumerable}s سټ: $\Obj u_0^n$، $\Obj u_1^n$، $\Obj u_2^n$، \dots
\end{enumerate}

په لومړۍ درجه منطق کښې عموماً د !!{identity}~$\eq$ نښه هم شامله وي۔ د ژبې~$\Lang{L}$ نامنطقي سمبولونه د مشخصو څيزونو، اړيکو او تابعو په باب د خبرو دپاره مهم دي۔ داسې !!{sentences} هم شته چې نامنطقي سمبول نۀ کاروي، خو يوازې له~$\eq$ سره د ژبې څرګندوونکے وس محدود وي۔ په دويمه درجه منطق کښې د اړيکو او تابعو د متغيرونو زېرمه بې‌پايه ده۔ نو د لومړۍ درجې په عبارت کښې د هر نامنطقي ثابت، محمول او تابعې سمبول پر ځاے هم‌ډوله ازاد متغير کارولے شو، چې د اړيکې يا تابعې د ځايونو شمېر ئې هماغه وي۔ متغيرونه داسې ټاکو چې کميت‌ښودونکے ئې ونۀ تړي او د متغيرونو ټاکنه هماغه څيز، اړيکه يا تابع ورکړي چې سمبول يې تعبيراوه۔ په دې ډول هماغه خبره د ځانګړو نامنطقي سمبولونو له زېرمې پرته څرګندېږي۔

\end{document}
